%!TEX root = ../template.tex
%%%%%%%%%%%%%%%%%%%%%%%%%%%%%%%%%%%%%%%%%%%%%%%%%%%%%%%%%%%%%%%%%%%%
%% abstract-en.tex
%%%%%%%%%%%%%%%%%%%%%%%%%%%%%%%%%%%%%%%%%%%%%%%%%%%%%%%%%%%%%%%%%%%%

\typeout{NT FILE abstract-en.tex}%

The Grande Panorama de Lisboa is a 23-metre azulejo tile panel attributed
to Gabriel del Barco (c.1700), the most complete surviving pictorial
record of Lisbon before the 1755 earthquake, catalogued at the closed
Museu Nacional do Azulejo (MNAz). Despite its scale, a single static
label leaves its 150+ depicted buildings uninterpretable without expert
knowledge. This report proposes TileStories, a
config-driven mobile augmented reality framework overlaying
visitor-adaptive content onto a heritage wall around a generalised,
wall-defined state axis --- instantiated on the Panorama as a four-epoch
timeline (pre-1755, the earthquake, the Pombaline reconstruction, and
today) --- with a short onboarding step adapting content depth and
routes to a small set of visitor profiles. The core technical
contribution is configuring and evaluating an existing
visual positioning service, Immersal, for a heritage-surface class ---
large, continuous, multi-point-of-interest tile walls --- via an
iterative Design Science Research approach with co-evolving problem
framing and technical exploration. The central technical risk is
demonstrated, not assumed: a working Immersal integration localises
against a real wall in two to three seconds. Generalisability is tested
across three structurally distinct cases: the Panorama; the
geometrically harder Chafariz Velho in Paço de Arcos, this work's
primary development surface given MNAz's closure; and a contemporary
wildlife mural with no equivalent status dimension. The completed thesis
is expected to contribute a validated multi-zone tracking configuration,
an adaptive visitor-profiling mechanism, a state-visualisation mechanism
generalised beyond a single wall, and an in-situ evaluation of
engagement, usability, and learning gain with real visitors.
 
\keywords{Augmented Reality, Cultural Heritage, Visitor Personalisation,
Mobile Applications, Historical Visualisation, Design Science Research,
Grande Panorama de Lisboa}
